
\documentclass[11pt,a4paper]{article}
\usepackage[hyperref]{naaclhlt2019}
\usepackage{times}
\usepackage{latexsym}
\usepackage{graphicx}

\usepackage{url}

\aclfinalcopy 

\title{Your proposed project title here}

\author{First Author \\
  {\tt email@domain} \\\And
  Second Author \\
  {\tt email@domain} \\\And
  Third Author \\
  {\tt email@domain} \\\And
  Fourth Author \\
  {\tt email@domain} \\}
  
% \setlength\textwidth{16.0cm}
\date{}

\begin{document}
\maketitle

\section{Introduction}
Tell us what problem you're going to work on. Provide some motivation for your idea: why is it interesting? Does it have any practical significance? 

Some general guidelines for this proposal document: 2 pages minimum for translation, 3 pages For ``choose your own projects'', due Mar 1. Please use this LaTeX template to write your proposal.

\section{Related work}

\noindent\textit{For both projects:} 
Check out papers at ACL/EMNLP/NAACL/TACL. Make sure to properly cite them. You can cite a paper parenthetically like this~\cite{andrew2007scalable} or use the citation as a proper noun, as in ``\newcite{borsch2011} show that...'' If you're not familiar with LaTeX, you'll have to add entries to \emph{yourbib.bib} to get them to show up when you cite them. 

\noindent\textit{For MT:} pick two papers on neural MT from ACL/EMNLP/NAACL/TACL and describe what they do. You may also choose papers from non-neural MT that you'd like to  extend to the neural setting.

\noindent\textit{For ``choose your own project'':}
Have others worked on this idea or related ideas? Clearly describe the some of these approaches, along with their pros and cons. You need to have at least five citations to related papers here. 

\section{Your approach}
How do you plan to solve the problem you chose? How will you approach it differently from previous work?

\noindent\textit{For ``choose your own project'':} remember that this project should take $\sim 2$ months of work! 

\paragraph{What baseline algorithms will you use?}
A baseline algorithm is one that is very simple and trivial to implement. For example, ``predict the most common class,'' or ``tag all capitalized words as names,'' or ``select the first sentence in the document''. Sometimes it can be difficult to get a fancy algorithm to beat a baseline. Always ask yourself, ``What's the simplest experiment I could do to (in)validate my hypothesis?'' Talented researchers have a knack for coming up with simple baselines.

\noindent\textit{For MT:} it could be a basic neural model (e.g., LSTM-LSTM), a non-neural phrase-based MT system, or something else. Be specific! 

\subsection{Milestones \& Schedule}
Divide your project into subtasks and estimate how much time each will take. If your group plans to divide subtasks amongst itself, also write who will be responsible for each milestone. If you plan to work on everything together, please say so here. Definitely budget some time for writing the progress report and final report, as well as performing an in-depth analysis of any models you build and/or data you collect. Sample schedule below:
\begin{enumerate}
    \item Acquire and preprocess data (1 week)
    \item Build models for task (3 weeks)
    \item Write progress report! (due Apr 1)
    \item Analyze the output of the model, do an error analysis (2 weeks)
    \item Work on final report and presentation (2 weeks)
\end{enumerate}

\begin{figure}[t]
    \centering
    \includegraphics[width=0.5\textwidth]{figs/sentence.png}
    \caption{Please feel free to include figures! If you want your figure to span both columns, use \emph{figure*} instead of \emph{figure}.}
    \label{fig:example}
\end{figure}

\section{Data}

\textit{For MT:}
Download the dataset (available on Piazza in the Resources section) and compute some basic statistics about it that you think will help inform your model design. Some obvious suggestions: number of sentence pairs, average number of tokens per sentence per language, size of vocab in both languages 

\noindent\emph{For ``choose your own project":} 
What text data do you plan to use in your project? Where will you get it from? Will you be annotating text yourselves? Convince us that it is available for you, and that you can easily get it, and that it is appropriate for the task and research questions you care about.

\section{Tools}
What existing libraries or toolkits are you going to use? Some questions to think about: will you be doing any preprocessing of your data such as tokenization or parsing? Will you be training logistic regression models? Will you be using deep learning libraries (if not, you need to justify why)? Will you need to use any services for GPUs?\footnote{As we said in class, we strongly suggest \url{https://colab.research.google.com}!} Do you need to use crowdsourcing?

\bibliographystyle{apalike}
\footnotesize
\bibliography{yourbib}


\end{document}
